(1) 利用镜像法，在虚构下半空间 \((0,0,-h)\) 放置电荷 \(-q\)，得到 \(\Phi=\frac{q}{4\pi\epsilon_0}[(\rho^2+(z-h)^2)^{-1/2}-(\rho^2+(z+h)^2)^{-1/2}]\)。代入 \(z=0\) 势恰为零，唯一性定理保证上半空间满足真实接地导体边值问题。
(2) 导体法向外电场为 \(E_z(0^+)=-\partial_z\Phi|_{0^+}=-qh/[2\pi\epsilon_0(\rho^2+h^2)^{3/2}]\)。因此表面感应电荷密度 \(\sigma(\rho)=\epsilon_0E_z(0^+)=-qh/[2\pi(\rho^2+h^2)^{3/2}]\)，负号表明正点电荷吸引电子。
(3) 面积分 \(Q_{\rm ind}=2\pi\int_0^\infty\sigma(\rho)\rho d\rho=-qh\int_0^\infty\rho(\rho^2+h^2)^{-3/2}d\rho=-q\)。镜像电荷距离真实电荷 \(2h\)，故导体对其施加 \(\mathbf F=-q^2\hat z/(16\pi\epsilon_0h^2)\)，方向向下。
(4) 由 \(F_z=-dU/dh\) 并选 \(U(\infty)=0\)，积分得 \(U(h)=-q^2/(16\pi\epsilon_0h)\)。若把镜像当作另一独立真实电荷来算相互作用能，会得到两倍负值；因为感应电荷分布随 \(h\) 改变，真实系统的机械势能必须由实际力积分求得。
